\begin{Solbook}{1}
 \par \noindent a) $-1200$; \hspace {2cm} b) 7; \hspace {2cm} c) 2022; \hspace {2cm} d) $-50$.  \hfill \qedEX 
\end{Solbook}
\begin{Solbook}{2}
 \par \noindent a) $x=28$; \hspace {1.5cm} b) $x=3$; \hspace {1.5cm} c) $x=2$; \hspace {1.5cm} d) $x=3$. hoặc $x=2$  \hfill \qedEX 
\end{Solbook}
\begin{Solbook}{3}
 \par \noindent 240 m. \hfill \qedEX 
\end{Solbook}
\begin{Solbook}{4}
 \par \noindent Cửa hàng A: $312\,000$ đồng và cửa hàng B: $301\,750$ đồng. Vậy nên mua ở cửa hàng B. \hfill \qedEX 
\end{Solbook}
\begin{Solbook}{5}
 \par \noindent 43 học sinh. \hfill \qedEX 
\end{Solbook}
\begin{Solbook}{6}
 \par \noindent 14 con chó, 22 con gà. \hfill \qedEX 
\end{Solbook}
\begin{Solbook}{7}
 \par \noindent  \begin {enumerate} \item Lớp 6E có nhiều học sinh nam nhất, lớp 6C có nhiều học sinh nữ nhất. \item Lớp 6E có sự chênh lệch nhiều nhất. \item Tổng số học sinh nam: 122, tổng số học sinh nữ: 98, tỷ lệ phần trăm: $55,45\%$. \end {enumerate}  \hfill \qedEX 
\end{Solbook}
\begin{Solbook}{8}
 \par \noindent a) A có 1007 số hạng, số hạng cuối cùng là: 4025. \qquad b) Có 702 trang. \hfill \qedEX 
\end{Solbook}
